\section{Arithmetic coefficients and an arbitrary original affine base}

The coefficient domain is a subring of the original characteristic-zero
field.  The affine base of the projective morphism is a separate algebra;
it may have zero divisors or nilpotents.

\begin{lemma}[The original affine base has a smooth arithmetic model]
\label{lem:kn-relative-arithmetic-affine-base-model}
\lean{EnvelopeProperty.Vanishing.RelativeArithmeticBaseModel.exists_original_base_model}
\uses{thm:arithmetic-original-finite-presentation-model,thm:arithmetic-smooth-base-same-model}
\leanok
Let $k$ be a field of characteristic zero and let $R$ be a finite-type
$k$-algebra.  Given a finite subset $E\subset k$, there are an integer
subalgebra $A\subset k$, smooth and of finite type over $\mathbb Z$,
containing $E$, a finitely presented $A$-algebra $B$, and a $k$-algebra
isomorphism $e:k\otimes_A B\simeq R$.  If
$\psi(b)=e(1\otimes b)$, then the square
\[
\begin{matrix}
\operatorname{Spec}R&\longrightarrow&\operatorname{Spec}B\\
\downarrow&&\downarrow\\
\operatorname{Spec}k&\longrightarrow&\operatorname{Spec}A
\end{matrix}
\]
with top map induced by $\psi$ is cartesian.  No domain or reducedness
condition is imposed on $R$ or $B$.
\end{lemma}
\begin{proof}
Since $k$ is noetherian, $R$ has a presentation
$k[T_1,\ldots,T_m]/(f_1,\ldots,f_s)$.  Let $A_0\subset k$ be the
integer algebra generated by $E$ and the coefficients of all the $f_i$,
and let $B_0=A_0[T_1,\ldots,T_m]/(f_1,\ldots,f_s)$.  Tensoring this
quotient presentation gives $k\otimes_{A_0}B_0\simeq R$, preserving the
original variables and coefficients.  The ring $A_0$ is an integer
domain of finite type.  Its generic point lies in its smooth locus
over $\mathbb Z$: its fraction field has characteristic zero and is
separably generated over $\mathbb Q$.  The smooth locus is open, so
some nonzero $a\in A_0$ has $A_0[1/a]$ smooth over $\mathbb Z$.
The injection into $k$ extends uniquely across this localization.
Identify its image with $A\subset k$ and put $B=A\otimes_{A_0}B_0$.
Finite presentation is preserved by this scalar extension, and tensor
associativity gives the required $e:k\otimes_A B\simeq R$.
The tensor product is the pushout of $A\to k$ and $A\to B$ in
commutative rings.  Transport its pushout square along $e$ and apply
the contravariant spectrum functor.  This gives exactly the displayed
cartesian square and its stated original field structure.
\end{proof}
